\phaseportada{PHASE 1}{Acquisition and Ground Truth}{
    % Tech Stack
    iOS Assistive Touch, PWA (JS/Netlify), Google Apps Script, Gemini Pro Vision.
}{
    % Core Challenge
    Market opacity and pervasive survivorship bias in standard platform exports.
}{
    % Key Insight
    Rejection data (the negative class) contains more predictive signal for policy modeling than acceptance data alone.
}{
    % Operational Outcome
    The successful deployment of a dual-engine telemetry system, yielding a consolidated ground-truth ledger of roughly 4,700 total market interactions.
}

\section{Phase 1: Acquisition and Ground Truth}

Project Pienza models the decision policy of an expert agent rather than a system in a learning phase. Prior to data collection, the subject had completed 24 months of field operations, concluding the \textit{Exploration} stage of the Reinforcement Learning cycle. 

Utilizing the framework of Hopfield Networks as explored by Anil Ananthaswamy in \textit{Why Machines Learn},\cite{ananthaswamy2024} the agent’s decision policy is modeled as a ``cooled'' system settled into a stable local minimum (\textit{an attractor state}). This equilibrium represents an optimization contingent on the agent’s specific constraints: geographic preferences, risk tolerance, and physical endurance. While a theoretical global optimum for city-wide revenue may exist in the state space, the captured dataset represents a phase of pure \textit{exploitation}. By prioritizing established, low-variance heuristics over high-energy experimental maneuvers, the resulting dataset provides a high-fidelity record of a converged policy.

\subsection*{Phase 1A: The 2-Engine Ingestion Infrastructure}
Project Pienza utilizes a proprietary, dual-engine acquisition ecosystem to overcome the data sparsity inherent in third-party platform exports. The primary acquisition campaign was executed over a strict 6-week observation window (August 22 -- October 1, 2025), digitizing the agent's operational reality in real-time.

\subsection*{Engine 1: The Operational Telemetry (GTS Webapp)}
To establish the ground truth of completed missions, a bespoke Progressive Web Application (PWA) designated as the \textbf{Geotimestamps (GTS) Webapp} was deployed. The interface was optimized for low-cognitive-load fieldwork, functioning as a ``One-Touch Timestamping'' instrument while maintaining resilience through a production-grade stack (Netlify Frontend, Google Sheets Backend).

\medskip
\noindent \textbf{Lifecycle Mapping Protocol (T0--T4).} The system logs the five critical state transitions of a gig economy mission with geospatial precision:
\begin{itemize}[noitemsep, topsep=3pt]
    \item \textbf{T0 (Search):} Idle state; actively seeking offers.
    \item \textbf{T1 (Acceptance):} Offer accepted; en route to pickup. \textit{(Captures quoted upfront fare)}.
    \item \textbf{T2 (Arrival):} At pickup location; waiting for passenger.
    \item \textbf{T3 (In-Progress):} Mission started.
    \item \textbf{T4 (Completion):} Mission finalized. \textit{(Captures realized net fare)}.
\end{itemize}

\medskip
\noindent \textbf{Technical Stack \& Evolution.} The system was engineered to maintain 100\% capture fidelity through two critical iterations:
\begin{itemize}[noitemsep, topsep=3pt]
    \item \textbf{Stateful Persistence (v4.0):} To prevent data loss during network intermittency, the app implements a \textit{Stateful LocalStorage Manager} and an \textit{Offline-First Queue System}. 
    \item \textbf{Asynchronous Geospatial Enrichment:} To overcome the geographic blindness of simple timestamps, a client-side pipeline orchestrates the \texttt{navigator.geolocation} API with reverse-geocoding services, enriching every state transition with high-precision coordinates prior to persistence.
\end{itemize}

\noindent \textbf{Output (\texttt{trip\_events}):} A verified timeline of physical state transitions and realized financial outcomes ($N \approx 350$ completed trips).

\subsection{Engine 2: The ``Total Offer'' Stream (OCR Pipeline)}
To neutralize \textbf{Survivorship Bias}, the system must capture the decision boundary—specifically the Rejected Offers (\textit{The Negative Class}). As market infrastructure provides no native access to historical offer logs, a custom ``Optical State Archival'' architecture was deployed.

The acquisition interface utilized an iOS device configured with an \textbf{Assistive Touch Macro}. This enabled a single-gesture capture protocol (The ``Manual Cookie'') that simultaneously acknowledged the offer cognitively and secured the raw visual data digitally. 

To ensure statistical representativeness, three operational mandates were enforced:
\begin{enumerate}
    \item \textbf{Full Spectrum Capture (Zero Filters):} All system-level destination filters and product category filters (Premium/Mid-tier/X) were disabled. This ensured the captured stream represented the unfiltered "True Market" liquidity, not a curated subset.
    \item \textbf{Safety-Driven Random Sampling:} Offers occurring during high-cognitive-load maneuvers or complex passenger interactions were intentionally skipped. This data loss is classified as \textbf{Missing Completely at Random (MCAR)}, introducing no systemic bias into the profitability models.
    \item \textbf{The "Driver-First" Learning Curve:} The dataset acknowledges a specific "Warm-Up" period (Week 1) where operational cadence took precedence over data capture, resulting in a small cluster of "Ghost Offers" (Completed missions without source image artifacts). These were programmatically handled in the Engineering Phase.
\end{enumerate}

The pipeline utilizes the \textbf{Google Gemini Pro Vision API} to transform visual assets into structured data. The primary challenge involved isolating the \textit{Foreground Signal} (the offer card) from \textit{Background Noise} (active navigation and status bar metadata). To ensure data integrity, the system prompt enforces strict scope limitations, instructing the model to ignore metadata outside the central offer card.

\noindent \textbf{Engine 2 Output (\texttt{offers}):} This high-fidelity dataset captures the total universe of opportunities, including the critical Negative Class ($N \approx 4,700$ total offers).

\noindent At this stage, both engines exist as separate, unlinked entities. In the subsequent engineering phase, these disparate records were unified into a SQL-queryable relational schema. 

\subsubsection{The Post-Session Reconciliation Protocol}
To maintain data integrity and preserve the agent's decision context, a rigorous post-session protocol was enforced immediately upon the completion of each fieldwork shift:

\begin{enumerate}
    \item \textbf{Telemetry Consolidation (GTS):} The raw, long-format event logs from the webapp were processed into a wide-format session ledger. This involved pruning duplicate entries and cancelled events, enabling the immediate calculation of session-level KPIs such as \textit{Net Spread} and \textit{Accumulated Deadhead}, metrics which would become engineered features in later stages.
    
    \item \textbf{High-Fidelity Cognitive Backtagging:} The Agent manually reviewed and tagged every rejected offer from that session to populate the multiclass target variable (\texttt{reason\_primary}). Executing this task same-day was imperative to capture the specific, contextual nuance of the decision before operational memory decay occurred.
\end{enumerate}


\subsection{Phase 1B: Data Staging and First Round of Feature Engineering}
Executed concurrently within the strict acquisition window, Phase 1B focuses on the immediate semantic structuring and target definition of the incoming datastream. 

\subsection{Core Entity: The \texttt{raw\_offers} Schema}

The ingestion pipeline consolidated OCR outputs into a primary staging table. The raw schema captured the unstructured state of the visual data prior to normalization.

\begin{table}[H]
    \centering
    \small
    \begin{tabular}{l l l}
        \toprule
        \textbf{Column} & \textbf{Data Type} & \textbf{Description} \\
        \midrule
        \texttt{ocr\_id} & Integer (PK) & Incremental sequential identifier assigned during ingestion \\
        \texttt{image\_filename} & String & Source asset filename stored in the archival vault. \\
        \texttt{time\_taken} & String & Raw time string (HH:MM) extracted from the device status bar. \\
        \texttt{ride\_type} & Categorical & Product class (e.g., \textit{X, mid-Tier, Premium}). \\
        \texttt{upfront\_fare} & Numeric & The nominal fare amount presented in the offer card. \\
        \texttt{pickup\_details} & String (Raw) & Combined text field containing ETA and distance to pickup. \\
        \texttt{pickup\_address} & String (Raw) & Unstructured pickup location text (street, POI). \\
        \texttt{trip\_details} & String (Raw) & Combined text field containing estimated trip duration and distance. \\
        \texttt{dropoff\_address} & String (Raw) & Unstructured destination location text. \\
        \texttt{rider\_rating} & String (Raw) & Combined text string with Passenger Rating and Trip Count. \\
        \texttt{special\_note} & Text (Blob) & Unstructured text field capturing all incentive tags (\textit{Surge, Turbo+}) and \\ 
        & & status flags (\textit{Reservation, Priority}). \\
        \bottomrule
    \end{tabular}
    \caption{Definition of the \texttt{raw\_offers} schema post-ingestion.}
    \label{tab:raw_schema}
\end{table}


\subsection{Target Engineering: \texttt{reason\_primary}}
The multiclass target variable, \texttt{reason\_primary}, is a hierarchical filtration system formalized as the \textbf{Cognitive Cascade}. To ensure a clean signal for supervised learning, each offer is assigned a single, mutually exclusive label representing the entry-point failure across a three-tiered triage: geospatial feasibility (Tier 1), economic viability (Tier 2), and strategic alignment (Tier 3). This prevents label contamination by focusing the model on the primary decision driver. The definitive taxonomy is outlined in Table \ref{tab:target_lexicon}.

\begin{table}[H]
    \centering
    \small
    \begin{tabularx}{\textwidth}{l X} 
        \toprule
        \textbf{Class Label} & \textbf{Operational Logic} \\
        \midrule
        \texttt{dropoff\_non\_operational} & Destination lies within a pre-defined zone outside the operational area. \\
        \addlinespace
        \texttt{dropoff\_proxy} & Destination is outside the primary zone but acceptable if aligned with a homecoming vector toward \textit{Anzures}. \\
        \addlinespace
        \texttt{low\_profitability} & Offer fails baseline EPH requirements relative to estimated duration. \\
        \addlinespace
        \texttt{long\_pickup\_time} & Uncompensated pickup time exceeds tolerance; threshold relaxes during extreme gridlock. \\
        \addlinespace
        \texttt{strategic\_mismatch} & High-value offer rejected due to unfavorable routing context (e.g., \textit{Santa Fe $\to$ Polanco} during Friday's peak gridlock). \\
        \addlinespace
        \texttt{expected\_value\_gamble} & Viable offer rejected based on the probabilistic expectation of a superior imminent event. \\
        \addlinespace
        \texttt{NULL} (Implicit) & Absence of objection signals an accepted offer. \\
        \bottomrule
    \end{tabularx}
    \caption{Lexicon of the Decision Target Variable (\texttt{reason\_primary}).}
    \label{tab:target_lexicon}
\end{table}








\subsection{Contextual Feature Extraction: \texttt{heuristic\_flag}}
The following table defines the categorical flags engineered to capture qualitative operational risks and strategic shifts that are not detectable through raw numerical features alone. These labels provide the model with the "Agent's perspective" on market conditions.

\begin{table}[H]
    \centering
    \small
    \begin{tabularx}{\textwidth}{l l X}
        \toprule
        \textbf{Category} & \textbf{Flag Label} & \textbf{Operational Logic} \\
        \midrule
        \multirow{3}{*}{\textbf{Risk Mitigation}} 
        & \texttt{deadhead\_risk} & Short-duration offers with a high probability of immediate return to an idle search state. \\
        & \texttt{long\_ride\_risk} & Trips $>45$ min where high traffic volatility creates risk of uncompensated time. \\
        & \texttt{dropoff\_uncertain} & Destinations with high geospatial entropy/ambiguity. \\
        \midrule
        \multirow{2}{*}{\textbf{Strategic Intent}} 
        & \texttt{obj\_end\_session} & Intent to terminate the work shift; prioritizes homecoming vector alignment over intrinsic EPH. \\
        & \texttt{friday\_gridlock} & Friday-specific regime; prioritizes egress from high-friction density sectors (\textit{Polanco}). \\
        \midrule
        \multirow{3}{*}{\textbf{System/Market}} 
        & \texttt{system\_error} & Physically impossible platform time-estimates (e.g., 20 min predicted time from Vistahermosa to Polanco during morning traffic). \\
        & \texttt{market\_anomaly} & Non-standard offer physics (e.g., operational outliers near National Holidays). \\
        & \texttt{protest\_anomaly} & External socio-political shocks (marches, road closures) disrupting standard market flow. \\
        \bottomrule
    \end{tabularx}
    \caption{Categorical Lexicon of Heuristic Flags for Model Enrichment.}
    \label{tab:heuristic_flags}
\end{table}












\subsection{Behavioral Targets and Operational Outcomes}
To differentiate between agent intent and realized performance, the schema implements two distinct categorical dimensions: \texttt{offer\_action} (Binary Intent) and \texttt{outcome} (Realized State).

\begin{table}[H]
    \centering
    \small
    \begin{tabularx}{\textwidth}{l l X} 
        \toprule
        \textbf{Target} & \textbf{Class Label} & \textbf{Strategic Context} \\
        \midrule
        \textbf{Intent} & \texttt{Accept / Reject} & Alternative binary target for supervised imitation learning. \\
        \midrule
        \textbf{Outcome} & \texttt{completed} & Transaction successfully finalized. \\
        & \texttt{rider\_canceled} & External termination by passenger post-acceptance. \\
        & \texttt{driver\_canceled} & Agent-initiated termination due to post-acceptance friction. \\
        & \texttt{system\_failure} & Offers intended for acceptance but lost due to operational capture latency. These are strategically mapped as \textbf{ACCEPTED} for model training to reflect the agent's true policy. \\
        & \texttt{NULL} & Implicit state for rejected offers; no operational outcome generated. \\
        \bottomrule
    \end{tabularx}
    \caption{Mapping of Behavioral Intent and Operational Outcome categories.}
    \label{tab:outcomes}
\end{table}









\subsection{Driver State Inference}
The variable \texttt{driver\_state\_at\_request} was engineered to classify the agent's operational context at the moment of an offer. This feature distinguishes between "Open" offers (received while \textit{Idle}) and "Chained" offers (received while \textit{On-Trip}), which is a primary driver of the agent's acceptance threshold.

The inference process followed an evolutionary path from manual verification to hybrid automation:

\begin{enumerate}
    \item \textbf{Verification Phase (Manual):} During the initial collection weeks, state assignment was performed via human-in-the-loop validation during the daily back-tagging protocol. The agent cross-referenced OCR timestamps against mission logs to establish a ground-truth baseline.
    
    \item \textbf{Automated Phase (SQL + Google Apps Script):} The final population of the state variable was executed through a two-step process. First, the ETL pipeline established the relational join between the \texttt{offers} and \texttt{trip\_events} tables. Second, a Google Apps Script (\texttt{code.gs}) was deployed to programmatically assign the \textit{On-Trip} state to any offer with a timestamp falling between a mission’s \texttt{T1} (Acceptance) and \texttt{T4} (Completion) boundaries.
\end{enumerate}






\subsection{Geospatial Enrichment}
To enable spatial clustering, the raw address strings required conversion to numerical coordinates. A serverless pipeline was architected using the \textbf{Google Geocoding API} within the Google Apps Script environment.

While the pipeline successfully populated the coordinate columns, a critical validation step was omitted: \textbf{Reverse Geocoding Verification}. The system naively accepted the API's top result as ground truth without auditing for semantic drift or "centroid defaults" (e.g., the API returning the center of a city for an ambiguous street name). This assumption of validity would later necessitate a massive "Deep Remediation" campaign during Phase 4, underscoring the risk of trusting black-box APIs without a human-in-the-loop audit.


\subsection{Acquisition Artifact: The Enriched \texttt{offers} Dataset}
By the conclusion of the acquisition campaign, the raw OCR stream had been enriched with manual behavioral labels, initial geospatial coordinates, and parsed operational metrics. The following schema represents the state of the dataset prior to the deep Data Engineering phase.

\begin{table}[H]
    \centering
    \small
    \begin{tabularx}{\textwidth}{l l X}
        \toprule
        \textbf{Feature Name} & \textbf{Data Type} & \textbf{Categories / Definition} \\
        \midrule
        \multicolumn{3}{c}{\textit{Source: OCR Extraction \& Initial Parsing}} \\
        \midrule
        \texttt{time\_taken} & String & Raw HH:MM timestamp from device. \\
        \texttt{product\_tier} & Categorical & \textit{X, mid-Tier, Premium}. \\
        \texttt{upfront\_fare} & Numeric & Nominal fare value (MXN). \\
        \texttt{special\_note} & String & Unstructured text (\textit{Surge, Turbo, Priority}). \\
        \texttt{star\_rating} & Numeric & Passenger rating (1.0--5.0 scale). \\
        \texttt{trip\_count} & Integer & Lifetime trips completed by passenger. \\
        \texttt{time\_to\_pickup} & Numeric & Parsed duration to pickup (seconds). \\
        \texttt{dist\_to\_pickup} & Numeric & Parsed distance to pickup (km). \\
        \texttt{est\_trip\_time} & Numeric & Parsed estimated trip duration (seconds). \\
        \texttt{est\_trip\_dist} & Numeric & Parsed estimated trip distance (km). \\
        
        \midrule
        \multicolumn{3}{c}{\textit{Source: Cognitive Backtagging \& Logic}} \\
        \midrule
        \texttt{offer\_action} & Binary & \textit{Accept, Reject}. \\
        \texttt{reason\_primary} & Categorical & \textit{dropoff\_non\_operational, low\_profitability, long\_pickup, strategic\_mismatch, expected\_value\_gamble}. \\
        \texttt{heuristic\_flag} & Categorical & \textit{deadhead\_risk, long\_ride\_risk, dropoff\_uncertain, obj\_end\_session, friday\_gridlock, system\_error, market\_anomaly, protest\_anomaly}. \\
        \texttt{driver\_state} & Categorical & \textit{Idle} (Open), \textit{On-Trip} (Chained). \\
        \texttt{outcome} & Categorical & \textit{Completed, Rider\_Cancel, Driver\_Cancel, System\_Failure}. \\
        
        \midrule
        \multicolumn{3}{c}{\textit{Source: Google Geocoding API}} \\
        \midrule
        \texttt{pickup\_lat/lon} & Float & Initial coordinate estimation. \\
        \texttt{dropoff\_lat/lon} & Float & Initial coordinate estimation. \\
        \bottomrule
    \end{tabularx}
    \caption{Final Schema State of the Acquisition Phase ($N \approx 4,700$).}
    \label{tab:acquisition_schema}
\end{table}
